% !TEX program = lualatex
% Transparencias de Fundamentos de las Matemáticas I
% Clase 03: Lógica proposicional

\documentclass[aspectratio=169,11pt]{beamer}

\usepackage{fontspec}
\usepackage[spanish,es-nodecimaldot]{babel}
\usepackage{amsmath,amssymb,mathtools}
\usepackage{booktabs}
\usepackage{csquotes}
\usepackage{xcolor}

\usetheme{Madrid}
\usecolortheme{default}
\setbeamertemplate{navigation symbols}{}
\setbeamertemplate{footline}[frame number]

\definecolor{FdMBlue}{RGB}{20,55,100}
\definecolor{FdMRed}{RGB}{130,30,30}
\definecolor{FdMGreen}{RGB}{25,100,60}

\setbeamercolor{structure}{fg=FdMBlue}
\setbeamercolor{title}{fg=white,bg=FdMBlue}
\setbeamercolor{frametitle}{fg=white,bg=FdMBlue}
\setbeamercolor{block title}{fg=white,bg=FdMBlue}
\setbeamercolor{block body}{bg=blue!3}
\setbeamercolor{alerted text}{fg=FdMRed}

\setmainfont{Latin Modern Roman}
\setsansfont{Latin Modern Sans}
\setmonofont{Latin Modern Mono}

\newcommand{\N}{\mathbb{N}}
\newcommand{\Z}{\mathbb{Z}}
\newcommand{\Q}{\mathbb{Q}}
\newcommand{\R}{\mathbb{R}}
\newcommand{\C}{\mathbb{C}}
\newcommand{\Pow}{\mathcal{P}}
\newcommand{\id}{\operatorname{id}}
\newcommand{\mcd}{\operatorname{mcd}}
\newcommand{\mcm}{\operatorname{mcm}}

\title[FdM I -- Clase 03]{Lógica proposicional}
\subtitle{Negación, conjunción, disyunción e implicación}
\author{Antonio Falcó Montesinos}
\institute{Universidad CEU Cardenal Herrera}
\date{Fundamentos de las Matemáticas I}

\begin{document}

\begin{frame}
  \titlepage
\end{frame}

\begin{frame}{Objetivos de la clase}
\begin{itemize}
  \item Reconocer proposiciones y conectores lógicos.
  \item Interpretar correctamente \(\neg\), \(\wedge\), \(\vee\), \(\Rightarrow\) y \(\Leftrightarrow\).
  \item Evitar errores con implicaciones y recíprocas.
\end{itemize}
\end{frame}

\begin{frame}{Mapa de la clase}
\tableofcontents
\end{frame}

\section{Proposiciones}

\begin{frame}{Proposiciones}
\begin{block}{Proposición}
\begin{itemize}
  \item Una proposición es una afirmación que puede ser verdadera o falsa.
  \item No toda expresión matemática es una proposición.
  \item El valor de verdad depende del contexto y de las definiciones previas.
\end{itemize}
\end{block}
\end{frame}

\begin{frame}{Proposiciones}
\begin{block}{Negación}
\begin{itemize}
  \item La negación de \(P\) se escribe \(\neg P\).
  \item Es verdadera exactamente cuando \(P\) es falsa.
  \item Negar no significa cambiar palabras al azar: hay que conservar el significado lógico.
\end{itemize}
\end{block}
\end{frame}

\section{Conectores}

\begin{frame}{Conectores}
\begin{block}{Conjunción y disyunción}
\begin{itemize}
  \item \(P\wedge Q\): se cumplen \(P\) y \(Q\).
  \item \(P\vee Q\): se cumple al menos una de las dos.
  \item En matemáticas, ``o'' suele ser inclusivo salvo indicación contraria.
\end{itemize}
\end{block}
\end{frame}

\begin{frame}{Conectores}
\begin{block}{Implicación}
\begin{itemize}
  \item \(P\Rightarrow Q\) se lee: si \(P\), entonces \(Q\).
  \item \(P\) es la hipótesis y \(Q\) la conclusión.
  \item Para refutar una implicación hace falta un caso con \(P\) verdadera y \(Q\) falsa.
\end{itemize}
\end{block}
\end{frame}

\begin{frame}{Conectores}
\begin{block}{Equivalencia}
\begin{itemize}
  \item \(P\Leftrightarrow Q\) significa que \(P\Rightarrow Q\) y \(Q\Rightarrow P\).
  \item Suele expresarse como ``si y solo si''.
  \item Demostrar una equivalencia exige probar las dos implicaciones.
\end{itemize}
\end{block}
\end{frame}

\section{Errores frecuentes}

\begin{frame}{Errores frecuentes}
\begin{block}{Recíproca y contrapositiva}
\begin{itemize}
  \item La recíproca de \(P\Rightarrow Q\) es \(Q\Rightarrow P\).
  \item La contrapositiva es \(\neg Q\Rightarrow \neg P\).
  \item La implicación y su contrapositiva son equivalentes; la recíproca no lo es en general.
\end{itemize}
\end{block}
\end{frame}

\section{Profundización}

\begin{frame}{Profundización}
\begin{block}{Tabla verbal de conectores}
\begin{itemize}
  \item \(P\wedge Q\): ambas afirmaciones deben cumplirse.
  \item \(P\vee Q\): basta con que se cumpla una de ellas.
  \item \(P\Rightarrow Q\): no afirma que \(P\) sea verdadera; afirma que \(Q\) se sigue de \(P\).
\end{itemize}
\end{block}
\end{frame}

\begin{frame}{Profundización}
\begin{block}{Implicación: caso crítico}
\begin{itemize}
  \item Una implicación \(P\Rightarrow Q\) solo falla cuando \(P\) es verdadera y \(Q\) es falsa.
  \item Si la hipótesis no se cumple, ese caso no refuta la implicación.
  \item Esta idea evita muchos errores al buscar contraejemplos.
\end{itemize}
\end{block}
\end{frame}

\begin{frame}{Profundización}
\begin{block}{Práctica rápida}
\begin{itemize}
  \item Escribir la recíproca de: si \(n\) es múltiplo de \(4\), entonces \(n\) es par.
  \item Escribir la contrapositiva del mismo enunciado.
  \item Decidir cuáles de las tres afirmaciones son equivalentes.
\end{itemize}
\end{block}
\end{frame}

\begin{frame}{Cierre}
\begin{itemize}
  \item La forma lógica determina el tipo de prueba.
  \item No se debe confundir una implicación con su recíproca.
  \item La próxima clase introduce cuantificadores.
\end{itemize}
\end{frame}

\begin{frame}{Trabajo recomendado}
\begin{itemize}
  \item Releer las secciones correspondientes del manual.
  \item Identificar definiciones, símbolos nuevos y enunciados principales.
  \item Preparar dudas y ejemplos para el seminario asociado.
\end{itemize}
\end{frame}

\end{document}
